\documentclass[12pt,a4paper]{article}

\usepackage[utf8]{inputenc}
\usepackage[T1]{fontenc}
\usepackage{amsmath,amssymb,amsfonts}
\usepackage{mathtools}
\usepackage{graphicx}
\usepackage[margin=2.5cm]{geometry}
\usepackage{setspace}
\usepackage{booktabs}
\usepackage{enumitem}
\usepackage[dvipsnames]{xcolor}
\usepackage{hyperref}
\usepackage{cleveref}
\usepackage{natbib}
\usepackage{abstract}
\usepackage{todonotes}
\usepackage{titlesec}
\usepackage{tikz}
\usetikzlibrary{shapes,arrows,positioning}

\hypersetup{
    colorlinks=true,
    linkcolor=blue!70!black,
    citecolor=blue!70!black,
    urlcolor=blue!70!black,
    pdfauthor={Judea Pearl},
    pdftitle={The Confounding of Scale: Why $\Delta_{13}$ Remains Unidentifiable},
}

\onehalfspacing
\titleformat{\section}{\large\bfseries}{\thesection.}{0.5em}{}
\titleformat{\subsection}{\normalsize\bfseries}{\thesubsection}{0.5em}{}
\renewcommand{\abstractnamefont}{\normalfont\bfseries}
\renewcommand{\abstracttextfont}{\normalfont\small}

\begin{document}

\begin{center}
    {\LARGE\bfseries The Confounding of Scale:\\[4pt]
    Why $\Delta_{13}$ Remains Unidentifiable}

    \vspace{1.2em}

    {\large Judea Pearl}\\[4pt]
    {\normalsize Cognitive Systems Laboratory, UCLA}\\
    {\small\texttt{judea@cs.ucla.edu}}

    \vspace{1em}
    {\small May 2026}
\end{center}
\vspace{0.5em}

\begin{abstract}
\noindent
Baldo \citep{baldo2026_scale} conjectures that the narrative residue ($\Delta_{13}$) will increase or remain constant with model scale, arguing that this would confirm "semantic gravity" over algorithmic failure. While the empirical prediction is testable, the proposed estimand ($\Delta_{13}$) remains causally unidentifiable. Because removing the narrative context $Z$ necessarily alters the prompt encoding $E$, the marginal difference $\Delta_{13}$ is hopelessly confounded. Measuring a confounded variable across different model scales does not un-confound it. I demonstrate via do-calculus that observing changes in $\Delta_{13}$ across scales cannot distinguish between an increase in causal injection (Mechanism C) versus an increase in encoding sensitivity (Mechanism B). To test the scale conjecture causally, one must measure the joint distribution $P(Y_A, Y_B \mid Z)$ across scales.
\end{abstract}

\section{The Unidentifiability of $\Delta_{13}$}

Baldo proposes measuring $\Delta_{13}$ across model scales to test whether "semantic gravity" is a fundamental law. The premise is that if $\Delta_{13}$ does not vanish, it must be because stronger semantic representations exert stronger narrative gravity over the logic.

However, as we established in previous causal evaluations, the intervention comparing Universe 1 to Universe 3 is confounded. Let $X$ be the combinatorial constraint graph, $Z$ be the narrative context, $E$ be the prompt tokenization encoding the graph, and $Y$ be the outcome.

The causal graph is: $Z \rightarrow E \rightarrow Y$, with an additional direct path $Z \rightarrow Y$ representing the hypothesized "semantic gravity" (Mechanism C).

The difference $\Delta_{13}$ attempts to measure the effect of intervening on $Z$: $P(Y \mid do(Z=1)) - P(Y \mid do(Z=0))$.

The fatal flaw is that we cannot intervene on $Z$ without also altering $E$. The prompt encoding $E$ acts as a backdoor path. Therefore, any measured difference $\Delta_{13}$ is an amalgamation of the true direct effect of the narrative (Mechanism C) and the spurious effect of altering the text encoding (Mechanism B). $\Delta_{13}$ cannot isolate either.

\section{Scaling a Confounded Variable}

Baldo suggests that if $\Delta_{13}$ increases with parameter scale, it confirms Mechanism C. This is a causal fallacy.

Let the model scale be $S$. As $S$ increases, the model's sensitivity to prompt encoding variations $P(Y \mid E)$ may also change. If a larger model is more sensitive to subtle changes in the tokenization of the constraint graph, then the spurious path $Z \rightarrow E \rightarrow Y$ will contribute a larger artifact to $\Delta_{13}$.

Observing $\Delta_{13} \ge \text{const}$ as $S \uparrow$ is perfectly consistent with a model that simply becomes more fragile to encoding shifts (Mechanism B), without any "causal injection" occurring.

\section{The Causally Valid Scale Test}

To properly test whether "semantic gravity" (Mechanism C) scales with model size, we must use an estimand that cleanly identifies Mechanism C. As Baldo has already conceded \citep{baldo2026_concession}, this requires the joint distribution test.

We must measure multiple independent boards within a shared narrative context and test for conditional independence: $P(Y_A, Y_B \mid Z) \neq P(Y_A \mid Z) P(Y_B \mid Z)$.

If the mutual information $I(Y_A; Y_B \mid Z)$ increases with model scale, then and only then can we conclude that the "semantic gravity" of the narrative frame is scaling. Measuring the confounded marginals $\Delta_{13}$ across scales will only yield uninterpretable noise.

\begin{thebibliography}{99}
\bibitem[Baldo(2026a)]{baldo2026_scale} Baldo, F.~S. (2026). The Scale Dependence Conjecture: Why Semantic Gravity is not a Transient Artifact. \emph{Procuradoria Geral do Estado de Rond\^onia}.
\bibitem[Baldo(2026b)]{baldo2026_concession} Baldo, F.~S. (2026). Mechanism C Identifiability: A Concession to Pearl and the Joint Distribution Test. \emph{Unpublished manuscript}.
\end{thebibliography}

\end{document}